\section{Summary}
\label{sec:summary}

In this work, we have proposed, simulated, and analyzed a Quantum Key Distribution (QKD) protocol based on the teleportation of a local charge observable. By studying two distinct Transverse Field Ising Models—one with a central star-coupling interaction ($H^{(1)}$) and another with nearest-neighbor interactions ($H^{(2)}$)—we have systematically compared the performance of charge teleportation against energy teleportation. Our investigation, conducted through both numerical simulations and Qiskit-based quantum circuit implementations, confirms that charge teleportation offers significant advantages for cryptographic applications.

Our primary finding is the superior stability and robustness of the charge-based protocol. While energy teleportation can, in certain regimes, yield a larger absolute signal ($\langle \Delta H_B \rangle$), it is highly sensitive to statistical noise and exhibits an asymmetric response to Alice's classical bit choice. In contrast, the teleported charge expectation value ($\langle \Delta Q_B \rangle$) shows a perfect mirror symmetry, which is crucial for a reliable bit-to-key mapping. This symmetry, combined with a signal that remains stable and scalable as the system size $N$ increases, makes the charge protocol inherently more resilient to the finite-shot noise that characterizes near-term quantum hardware.

We also characterized the protocol's dependence on the system's physical parameters. The star-interaction Hamiltonian provides a stronger teleportation signal due to direct Alice-Bob coupling, whereas the signal in the more realistic nearest-neighbor model decays with distance. For the latter, the choice of Alice's measurement basis ($\sigma_A = X_0$ or $Y_0$) was shown to be a critical parameter influencing protocol efficiency. Critically, our protocol was validated on real quantum hardware for a minimal two-qubit system ($N=1$). Despite the presence of device noise, the experimental results successfully reproduced the characteristic separation of Bob's observable expectation values, confirming the fundamental viability of the QKD mechanism. Finally, our cryptographic security analysis quantified the precise trade-off between physical noise and secure key rate, demonstrating that the protocol is secure for $N=2$ (up to a noise threshold of $p_{\rm th} \approx 0.02$) but that longer chains ($N \ge 3$) would require further optimization of the Hamiltonian parameters to overcome high intrinsic error rates.

Looking forward, this work opens several avenues for future research. One promising direction is to systematically map other many-body Hamiltonians to their global symmetries to identify a wider set of conserved charges suitable for teleportation. While this paper demonstrated the utility of one such charge operator, a broader library could enable the tailoring of the protocol to specific physical systems for optimal performance. Furthermore, recent developments in \textit{Timelike} Quantum Energy Teleportation (TQET) offer an intriguing path to enhancing signal strength \cite{Ikeda2025_TQET}. TQET utilizes temporal correlations by introducing a waiting period before Bob's operation, which has been shown to boost energy transfer efficiency by over an order of magnitude. Since QKD protocols are often valued for security over speed, incorporating a timelike approach could dramatically improve the signal-to-noise ratio of charge teleportation, making it even more robust for real-world applications. The integration of these concepts could lead to a new generation of highly practical and secure QKD systems built upon the versatile framework of observable teleportation.
